\begin{table}[H]
	\caption{\textit{Non-Functional Requirement}}
\label{tab:kebutuhan nonfungsional}
\begin{tabular}{ | p{2cm} | p{3cm} | p{7cm} |}
	\hline
	\textbf{Kode} & \textbf{Aspek} & \textbf{Nama Kebutuhan} \\ \hline
	NF1 & \textit{Performance} & Sistem dapat memproses rata-rata satu foto dalam waktu kurang dari 1 detik menggunakan spesifikasi perangkat keras standar \\ \hline
	NF2	& \textit{Performance} & Sistem memiliki tingkat akurasi (\textit{accuracy}) minimal 85\% dalam mengidentifikasi foto pribadi \\ \hline
	NF3	& \textit{Reliability} & Sistem tidak mengalami \textit{crash} (berhenti mendadak) saat menemui \textit{file} gambar yang rusak (\textit{corrupt}), format tidak didukung, atau foto tanpa wajah \\ \hline
	NF4	& \textit{Scalability} & Sistem harus mampu memproses masukan \textit{dataset} dalam jumlah besar (minimal 500 - 1.000 foto) dalam satu kali eksekusi tanpa mengalami kegagalan memori \\ \hline
	\end{tabular}
\end{table}